% C-counts.tex: the numbers, and how each was measured.
\section{Counts at 1.0.0-rc.7, and how they were measured}
\label{app:counts}

Every count in this paper was taken from the tree at version \code{1.0.0-rc.7} (commit
\code{4e74376}, 20 September 2026) by the command beside it, not copied from an earlier document.
Where the repository's own stamped number was measured at an earlier version, the version is
stated. Where a number is external, the row says whose it is.

Version~3 of this paper was written at \code{1.0.0-rc.1} and re-measured here. Fourteen of the
twenty-six quantities below are unchanged; the rows that moved are marked in the text where they
matter. One method changed rather than one number: the route count is the same and the command
that produced it no longer is.

{\footnotesize
\setlength{\tabcolsep}{4pt}\renewcommand{\arraystretch}{1.12}
\rowcolors{2}{tabstripe}{white}
\begin{xltabular}{\textwidth}{>{\RaggedRight\arraybackslash\hspace{0pt}}p{2.1in} >{\raggedleft\arraybackslash}p{0.6in} >{\RaggedRight\arraybackslash\hspace{0pt}}X}
\caption{Counts at \code{1.0.0-rc.7} (commit \code{4e74376}) and their method. The tree moves after the measurement; a count taken later can differ, and the commit is what makes each one reproducible.}\label{tab:counts}\\
\toprule
\rowcolor{tabheader}\thead{Quantity} & \thead{Count} & \thead{Method} \\
\midrule
\endfirsthead
\toprule
\rowcolor{tabheader}\thead{Quantity} & \thead{Count} & \thead{Method} \\
\midrule
\endhead
\midrule\multicolumn{3}{r}{\itshape continued on the next page}\\
\endfoot
\bottomrule
\endlastfoot
Canonical tables in \code{01\_schema.sql} & 45 & The file declares 50 \code{CREATE TABLE} statements, of which five are the default partition children of the partitioned event tables. A migrated deployment holds 52 tables. \\
Triggers in \code{06\_triggers.sql} & 42 & \code{grep -c 'CREATE TRIGGER'} \\
Triggers in the loaded schema & 46 & The file's 42, and four declared beside what they guard: one in the migration registry's file, one in \code{01\_schema.sql}, two in \code{12\_v7\_constraints.sql}. Read from \code{pg\_trigger} on a loaded and migrated database, excluding internal triggers and the 25 copies PostgreSQL makes on the partitions of partitioned tables \\
Tables under an append-only trigger & 31 & The \code{BEFORE UPDATE OR DELETE} triggers of \code{06\_triggers.sql} and the migrations, by table; the design record names fourteen audit-of-record instances among them \\
Invariant checks & 313 & \code{grep -c '\^{}def check\_' polaris\_checks/checks.py}, equal to the registered list; each has a detection test with a positive control in \code{test\_checks.py} \\
Mutation drills & 12 & Files matching \code{scripts/polaris-*mutation-drill*}: checks, constraints, triggers, procedures, the zero-knowledge witnesses, the conformance contract, the two SDKs, the OpenID4VP verifier, the application, coverage, the constitution, and moves between modules \\
Application routes & 124 & \code{grep -c '\^{}@app.route('} across \code{polaris\_web/*.py}. The count is unchanged from \code{1.0.0-rc.1}; the METHOD is not. On 18 September 2026 \code{app.py} was decomposed into ten route modules, and the old method, which read that one file, now returns 13 \\
CI jobs in \code{ci.yml} & 20 & The job keys of the workflow; two more workflows run monthly and weekly, and a manual one publishes \\
Signed wire formats in the specification & 27 & Distinct \code{polaris-*/1} strings in \code{docs/reference/WIRE-SPEC.md} \\
Conformance cases & 118 & The cases of \code{conformance/cases.json}, across twenty-five artifact types; the frozen version-1 set holds 44 \\
Verdict fields no published case constrains & 44 & \code{SURVIVORS\_EXPECTED} of \code{scripts/polaris-conformance-mutation-drill.py}, which the drill holds to its own measurement in both directions: of 89 decision fields across the 27 verifiers the cases reach, 44 survive being forced permissive, and none of them could be constrained by writing a case \\
Formal models checked in CI & 4 & \code{meta/tla/*.tla}; three carry a counterpart configuration they must fail \\
Drill scripts & 80 & Files matching \code{scripts/polaris-*drill.*}, the twelve mutation drills among them \\
Migrations & 45 & Files matching \code{polaris\_sql/migrations/*.up.sql}, each with a revert \\
Operator runbooks and ledgers & 21 & Markdown files in \code{docs/operator/}, the index included \\
Tagged versions & 3 & \code{git tag}. At \code{1.0.0-rc.1} this was 295, one per ship from \code{v9.30} (16 May 2026) to \code{v9.466} (15 September 2026). A tag marks a RELEASE and not a ship, so those were removed; what remains is \code{v1.0.0-rc.1}, \code{-rc.2} and \code{-rc.3}, the three published candidates \\
Test functions defined & 2,189 & \code{grep -c 'def test\_'} across the product, CLI, check-layer, script, card, simulation, SDK and package suites; a definition count, not a run count \\
Product tests passing, as stamped by the README & 1,030 & Measured at \code{1.0.0-rc.7} on the reference machine, \code{pytest -q} per suite; three skip without real ML-DSA \\
Crypto witness tests, as stamped & 107 of 112 & Measured at \code{1.0.0-rc.7}, unchanged from \code{1.0.0-rc.1}; five need a PKCS\#11 module or a real KMS key, and CI runs all but the KMS one \\
OpenID4VP verifier package tests & 244 & \code{pytest} in \code{packages/polaris-oid4vp}; three are the positive control without which the refusals are vacuous \\
Packages on public registries & 4 & \code{polaris-verify}, \code{polaris-oid4vp}, \code{polaris-sdk-python} on PyPI; \code{polaris-sdk-ts} on npm; each verified by installing from the live registry. All four served \code{1.0.0-rc.1} when Version~3 was written. \code{1.0.0-rc.3}, published 18 September 2026, is a pre-release: an installer takes it only when asked (\code{pip install --pre}, the npm tag \code{next}), and a plain install still resolves \code{0.1.0}, the last stable release. Read off the live registries on 23 September 2026 \\
Hosted conformance modules & 11 & The OpenID Foundation's service, 15 September 2026: 7 \code{PASSED}, 4 \code{REVIEW}, 0 failures, 0 warnings; the service's own signed exports \\
External test vectors for ML-DSA-65 & 58 of 58 & Project Wycheproof, pinned to a commit, under both witnesses on every push \\
Named outside implementations that have presented & 1 & walt.id Wallet API v2 at a pinned image digest, 15 September 2026; one path, one format \\
Fixes caused by an external finding & 2 & The scoreboard's row: the certificate the wallet refused (15 September 2026), and a version stamp an outside reviewer could not read unambiguously (17 September 2026) \\
Independent security reviews, pilots, non-author operators & 0 & The scoreboard's rows, blank \\
\end{xltabular}
}

The paper's performance numbers are those of Table~\ref{tab:perf}, each with the version that
measured it; the dyno and the application baseline have not been re-run since \code{v9.278} and
\code{v9.191} and are cited with those stamps. The reference machine throughout the repository is an
Apple Silicon laptop with eight cores and sixteen gigabytes of memory.
